\documentclass[11pt]{amsart}
\usepackage[T1]{fontenc}
\usepackage{lmodern}
\usepackage{amsmath,amssymb,amsthm}
\usepackage[margin=1.1in]{geometry}
\usepackage{booktabs}
\usepackage{hyperref}

\newtheorem{theorem}{Theorem}
\newtheorem{lemma}[theorem]{Lemma}
\newtheorem{proposition}[theorem]{Proposition}
\newtheorem{corollary}[theorem]{Corollary}
\newtheorem{conjecture}[theorem]{Conjecture}
\newtheorem{question}[theorem]{Question}
\theoremstyle{definition}
\newtheorem{definition}[theorem]{Definition}
\newtheorem{remark}[theorem]{Remark}
\newtheorem{example}[theorem]{Example}

\newcommand{\Ev}{\mathrm{Ev}}
\newcommand{\Od}{\mathrm{Od}}
\newcommand{\cR}{\mathcal R}
\newcommand{\cD}{\mathcal D}

\title[Route descent]{Route descent: transversal realisability, minimum $T$-joins,\\ and Lov\'asz
decompositions that separate vertices}
\author{Justin Williams}
\date{November 2026 (working draft)}

\begin{document}
\begin{abstract}
Let $G$ be a connected graph of even order $n$. The cut-selection theorem of the endpoint min-max
manuscript pairs the even-degree vertices of $G$ along edge-disjoint \emph{routes} in a forest $T$-join, so
that every \emph{transversal} (one vertex from each pair) satisfies all cut conditions for a partition of
$E(G)$ into $n/2$ paths with two ends at each chosen vertex. The conjecture that every transversal is
realisable (TRC) implies Gallai's conjecture. We study the natural induction for TRC, which deletes
one route and puts it back.
\begin{itemize}
\item \emph{Route descent.} The class of route systems is closed under deleting a route, its base case is
Lov\'asz's theorem, and the inductive step is equivalent to finding a realisation in which the deleted
route is a terminal segment of a path. We show that TRC for some root-containing transversal implies
every level of the branching hypothesis of the Lean spine, hence Gallai.
\item \emph{Minimum $T$-joins.} If the $T$-join is minimum, every segment of every route is a geodesic,
minimality survives route deletion, and the only obstruction ever met with spanning-tree $T$-joins
(an unmarked end of degree two whose route runs through its other neighbour) cannot occur. For any
route system, a realisation in which \emph{all} routes are terminal segments is equivalent to a
\emph{constrained Lov\'asz decomposition} of the free graph, in which every vertex is the end of exactly
one path and the paths at the unmarked ends avoid their routes. We conjecture that such decompositions
exist for minimum $T$-joins.
\item \emph{Separation.} We prove a strengthening of Lov\'asz's theorem. If a graph has at most one vertex
of positive even degree, and $X$ is a set of vertices at pairwise distance at least three, then it has a
path decomposition with the parity-forced ends in which every path meets $X$ at most once. In
particular any two vertices at distance at least three can be separated. The proof runs
Lov\'asz's induction and uses that Fan's addible-edge lemma changes the vertex set of only one path, by
one vertex. We show by example that distance three cannot be replaced by two.
\end{itemize}
We also prove TRC for two classes reachable by Fan's lemma, state the constrained Lov\'asz conjectures that
would give TRC for minimum $T$-joins, and record what we know of their status. One of them would imply
Gallai's conjecture for every graph whose even vertices induce a subgraph with a perfect matching. The
proofs use no computation. Computations (an exact oracle, about $1.5\times10^5$ transversals and
$6\times10^5$ extension queries) are reported separately and only guided the statements. The
separation theorem, the vertex-set form of Fan's lemma and the reduction of Gallai's conjecture to an
existential TRC are formalised in Lean.
\end{abstract}
\maketitle

\section{Introduction}\label{sec:intro}

All graphs are finite and simple. A \emph{path partition} (or decomposition) of a graph $G$ is a partition
of $E(G)$ into paths, each with at least one edge. Gallai conjectured that every connected graph on $n$
vertices has a path partition into at most $\lceil n/2\rceil$ paths. Lov\'asz~\cite{L68} proved this when
at most one vertex has even degree. Pyber~\cite{Pyb96} proved it when the even vertices induce a forest,
Fan~\cite{Fan05} when every block of that subgraph is triangle-free of maximum degree at most three,
and Botler and Sambinelli~\cite{BS21} extended Fan's theorem. The conjecture is also known for maximum
degree at most five~\cite{BP19}, among other classes.

This manuscript belongs to a programme that reduces Gallai's conjecture to statements about where path
ends can be placed. For a graph $G$ of even order $n=2k$, write $\Ev$ and $\Od$ for its sets of
vertices of positive even degree and of odd degree. For $S\subseteq\Ev$ the \emph{simple profile} $b_S$ is
$1$ on $\Od$, $2$ on $S$ and $0$ on $\Ev\setminus S$, and $S$ is \emph{realisable} if $E(G)$ has a path
partition with exactly $b_S(v)$ ends at every vertex $v$. Such a partition has $|\Od|/2+|S|$ paths. So a
realisable $S$ with $|S|=|\Ev|/2$ gives Gallai's bound.

The cut-selection theorem~\cite[Theorem~21]{EMM} constructs a pairing $P$ of $\Ev$ by edge-disjoint
paths, the \emph{routes}, inside a forest $T$-join for $T=\Ev$. Every transversal of $P$ then satisfies
all the cut inequalities necessary for realisability. Since cut inequalities are not sufficient, and
since realisability of a prescribed half-set is NP-complete~\cite[\S\S5--6]{EMM}, the right statement
concerns the transversals of these pairings only:

\begin{conjecture}[TRC, transversal realisability~\cite{EMM}]\label{conj:trc}
For every connected graph $G$ of even order, every pairing $P$ produced by the cut-selection construction,
and every transversal $M$ of $P$, the set $M$ is realisable.
\end{conjecture}

TRC implies Gallai's conjecture (Proposition~\ref{prop:spine}). This manuscript studies the induction
that deletes one route and puts it back, which we call \emph{route descent}.

\subsection*{Results}
\begin{enumerate}
\item Route descent is well founded: the class is closed under deleting a route, the base case is
Lov\'asz's theorem, and the inductive step at a route $R$ succeeds exactly when some realisation has $R$
as a terminal segment of a path (Section~\ref{sec:descent}). A realisable transversal containing a
prescribed vertex gives all levels of the branching hypothesis $\mathrm{BranchBase}_k$ of the Lean spine
at once (Proposition~\ref{prop:spine}).
\item For minimum $T$-joins, route segments are geodesics, minimality is hereditary, and the
\emph{trivial obstruction} (the unmarked end has no edge leaving its route) is impossible
(Section~\ref{sec:min}).
\item Realisations in which every route is terminal are exactly the \emph{constrained Lov\'asz
decompositions} of the free graph $G-J$ (Proposition~\ref{prop:rr}). This turns TRC into a problem about
graphs with all degrees odd.
\item \emph{Separation theorem} (Theorem~\ref{thm:sep}): for every set $X$ of vertices at pairwise
distance at least three, some Lov\'asz decomposition has every path meeting $X$ at most once; in
particular, it can separate any two vertices at distance at least three. The proof rests on a vertex-set form of Fan's lemma
(Lemma~\ref{lem:fan}), proved in full. Both are formalised in Lean (Section~\ref{sec:lean}).
\item TRC holds for \emph{Fan-isolated} route systems, and every orientation of a unique perfect matching
of $G[\Ev]$ is realisable (Section~\ref{sec:classes}). Both are proved with Fan's lemma alone.
\end{enumerate}
Section~\ref{sec:conj} states the conjectures that would complete the programme for minimum $T$-joins,
and records their computational status. Section~\ref{sec:compare} compares the results with the known
classes.

\section{Preliminaries}\label{sec:prelim}

For a path partition $\cD$ of a graph $K$ and a vertex $v$, $b_\cD(v)$ is the number of paths of $\cD$
having $v$ as an end; a path has two distinct ends. Then $b_\cD(v)\equiv d_K(v)\pmod 2$, since each path
through $v$ uses two edges at $v$ and each path ending at $v$ uses one. We call $\cD$ a \emph{Lov\'asz
decomposition} of $K$ if $b_\cD(v)=d_K(v)\bmod 2$ for every vertex $v$. So every odd vertex is the end of
exactly one path, and every even vertex is the end of none. For a vertex $c$ of odd degree, $Q_c$ denotes
the unique path of a Lov\'asz decomposition having $c$ as an end.

\begin{theorem}[Lov\'asz~\cite{L68}, strong form]\label{thm:lovasz}
If at most one vertex of $K$ has positive even degree, then $K$ has a Lov\'asz decomposition.
\end{theorem}

This form is proved in the Lean library (\texttt{lovasz\_ends}); the proof of
Theorem~\ref{thm:sep} below reproves it.

For vertices $t\neq u$ with $tu\notin E(K)$ we write $K+tu$ for the graph with the edge $tu$ added. For a
path partition $\cD$ of $K$ let $Z_\cD(u)=\{w\in N_K(u): b_\cD(w)=0\}$.

\begin{lemma}[Fan's lemma with vertex sets; cf.~\cite[Lemma~3.4]{Fan05}]\label{lem:fan}
Let $\cD$ be a path partition of $K$, and let $t\neq u$ be vertices with $tu\notin E(K)$ and
$b_\cD(t)>|Z_\cD(u)|$. Then $K+tu$ has a path partition $\cD^+$ with $|\cD^+|=|\cD|$,
$b_{\cD^+}=b_\cD-\mathbf 1_t+\mathbf 1_u$, and a bijection $\varphi:\cD\to\cD^+$ with $V(\varphi(Q))=V(Q)$
for all $Q\in\cD$ except one, $Q_0$, for which $V(\varphi(Q_0))=V(Q_0)\cup\{u\}$.
\end{lemma}

\begin{proof}
For a vertex $x\neq u$ and a path $Q\in\cD$ having $x$ as an end and containing $u$, read $Q$ from $x$ and
let $\pi(x,Q)$ be the vertex just before $u$; it is a neighbour of $u$ different from $u$. A \emph{chain}
is a sequence of distinct vertices $t=x_0,x_1,\dots,x_k$ together with paths $Q_0,\dots,Q_{k-1}\in\cD$ such
that $Q_i$ has $x_i$ as an end, contains $u$, and $\pi(x_i,Q_i)=x_{i+1}$, and such that some path
$Q^\ast\in\cD$ having $x_k$ as an end does not contain $u$.

\emph{A chain exists.} Let $X$ be the set of vertices reachable from $t$ by steps $x\mapsto\pi(x,Q)$. Then
$u\notin X$, because $t\neq u$ and $\pi$ never takes the value $u$. A shortest sequence of steps from $t$ to
a vertex of $X$ consists of distinct vertices. So if no chain exists, then every path having an end in $X$
contains $u$. Map each pair $(x,Q)$ with $x\in X$ and $Q$ a path ending at $x$ to the edge
$\pi(x,Q)u$ of $Q$. This map is injective: edges of distinct paths are distinct, and if $Q$ has both
ends in $X$, reading $Q$ from its two ends meets the two different neighbours of $u$ on $Q$. Its image
consists of edges $yu$ with $y\in X\cap N_K(u)$. Hence
\[
\sum_{x\in X}b_\cD(x)\le|X\cap N_K(u)|.
\]
On the other hand $t\in X\setminus N_K(u)$, and every vertex of $X\cap N_K(u)$ outside $Z_\cD(u)$ has at
least one end, so the left side is at least $b_\cD(t)+|X\cap N_K(u)|-|Z_\cD(u)|$. Thus
$b_\cD(t)\le|Z_\cD(u)|$, a contradiction.

\emph{Executing a chain.} Write $K_i=K+x_iu-x_{i+1}u$ for $i<k$ and $K_k=K+tu$. We construct partitions
$\cD_0=\cD$ of $K$ and $\cD_{i+1}$ of $K_i$ as follows.

\emph{Re-routing} at step $i<k$: the path $Q_i$, read from $x_i$, is $x_i\cdots x_{i+1}\,u\cdots y$.
Replace it by $x_{i+1}\cdots x_i\,u\cdots y$, that is, reverse the segment from $x_i$ to $x_{i+1}$, and
trade the edge $x_{i+1}u$ for the edge $x_iu$. This is again a path, with the \emph{same vertex set}. Its
ends are $x_{i+1}$ and $y$. So the end at $x_i$ moves to $x_{i+1}$, and the edge $x_{i+1}u$ becomes free.
Read from its other end $y$, the path is unchanged up to $u$.

\emph{Appending} at step $k$: replace $Q^\ast$ by $Q^\ast+x_ku$, which is a path, since $u\notin V(Q^\ast)$.
Its vertex set gains exactly $u$, and the end at $x_k$ moves to $u$.

It remains to check that each step acts on the intended path. The paths $Q_i$ are pairwise distinct
except possibly when $Q_i=Q_j$ with $j<i$. In that case $x_i$ is the other end of $Q_j$, since
$x_i\neq x_j$. The re-routing at step $j$ left $Q_j$ unchanged when read from that end up to $u$, so
$\pi(x_i,\cdot)$ is still $x_{i+1}$, and step $i$ applies to the current form of $Q_j$. A path cannot be
used three times, because it has only two ends and the $x_i$ are distinct. The path $Q^\ast$ avoids $u$,
so it is none of the $Q_i$, and it is untouched until step $k$. Also, at step $i$ the edge $x_iu$ is
absent from $K_{i-1}$ and the edge $x_{i+1}u$ lies on $Q_i$, so every step is legal. Finally, the ends
change only at $t$, which loses one, and at $u$, which gains one; every $x_i$ with $0<i\le k$ gains one
end at step $i-1$ and loses it at step $i$. The vertex sets change only at step $k$.
\end{proof}

This is Fan's proof, formalised in the Lean library as \texttt{fan\_exists\_chain},
\texttt{reroute\_move}, \texttt{append\_move} and \texttt{fan\_exec}. The statement of
Lemma~\ref{lem:fan} records the extra information that the vertex sets change only at the final
append. This observation is the engine of Theorem~\ref{thm:sep}.

\section{Route systems and route descent}\label{sec:descent}

Throughout, $G$ is a graph without isolated vertices whose components have even order. Then $|\Ev|$ is
even (it is $n$ minus the even number $|\Od|$).

\begin{definition}\label{def:routes}
A \emph{route system} for $G$ is a pairing $P$ of $\Ev$ together with pairwise edge-disjoint paths
$R_p$ ($p\in P$) of $G$, where $R_p$ joins the two vertices of $p$. Its \emph{union} is
$J=\bigcup_pE(R_p)$, and $F=G-J$ is the \emph{free graph}. A \emph{passage} at $v$ is a pair $\{x,y\}$ such
that $x,v,y$ are consecutive on a route, and if $v\in\Ev$ the \emph{end edge} at $v$ is the edge of the
route ending at $v$. For a vertex $v$, $F^v$ is $F$ with the edges at $v$ deleted and $v$ joined to every
vertex of $V-v-N_G(v)$. The route system is a \emph{Thm-21 system} if $J$ is a forest and, for every $v$,
there are pairwise edge-disjoint paths of $F^v$, one joining $x$ to $y$ for each passage $\{x,y\}$ at $v$,
and one joining $z$ to $v$ for the end edge $vz$ if $v\in\Ev$ (\emph{linked transitions}).
\end{definition}

Theorem~21 of~\cite{EMM} produces a Thm-21 system on every connected graph of even order, from any
forest $T$-join, and every transversal of its pairing is cut-feasible. TRC asserts that these
transversals are realisable.

\begin{lemma}[basic facts]\label{lem:basic}
For a route system, $J$ is a $T$-join for $T=\Ev$, and every vertex has odd degree in $F$. An interior
vertex of a route has degree at least three.
\end{lemma}

\begin{proof}
A vertex of $\Ev$ is the end of exactly one route and may be interior to others; an odd vertex is interior
only. So $d_J(v)$ is odd exactly on $\Ev$, and $d_F(v)=d_G(v)-d_J(v)$ is odd for every $v$. An interior
vertex $v$ has $d_J(v)\ge2$. If $d_G(v)=2$, then $v\in\Ev$, so $v$ is also an end of a route and
$d_J(v)\ge3>d_G(v)$, which is impossible.
\end{proof}

\begin{lemma}[upward closure]\label{lem:up}
If $S\subseteq\Ev$ is realisable and $x\in\Ev\setminus S$, then $S\cup\{x\}$ is realisable.
\end{lemma}

\begin{proof}
In a realisation of $S$ the vertex $x$ has no end and positive degree, so it is interior to some path.
Splitting that path at $x$ gives two paths with an end at $x$ each, and changes nothing else.
\end{proof}

Hence TRC is equivalent to: every $S\subseteq\Ev$ meeting every pair of $P$ is realisable.

\begin{lemma}[heredity]\label{lem:hered}
Let $(P,\cR)$ be a route system for $G$, let $p=\{a,c\}\in P$, and let $G'=G-E(R_p)$. Then $G'$ has no
isolated vertex, every component of $G'$ has even order, $\Ev(G')=\Ev\setminus\{a,c\}$, and
$(P-p,\cR-R_p)$ is a route system for $G'$ with union $J-E(R_p)$ and the same free graph $F$. If
$(P,\cR)$ is a Thm-21 system, so is $(P-p,\cR-R_p)$.
\end{lemma}

\begin{proof}
Deleting $E(R_p)$ lowers the degrees of the interior vertices of $R_p$ by two and those of $a,c$ by one.
So parities change only at $a$ and $c$, which become odd. Since $a,c\in\Ev$ have degree at least two, they
keep positive degree. An interior vertex keeps degree at least one, by Lemma~\ref{lem:basic}. Hence
$\Ev(G')=\Ev\setminus\{a,c\}$, which is paired by $P-p$ through the remaining routes, and $J'=J-E(R_p)$
satisfies $G'-J'=G-J$. Each component $C$ of $G'$ contains an even number of odd vertices of $G'$, by the
handshake lemma, and an even number of vertices of $\Ev(G')$, since routes lie inside components and pair
them. So $|C|$ is even.

For the Thm-21 property: $J'$ is a forest. For each vertex $v$ we have $N_{G'}(v)\subseteq N_G(v)$, so
$F'^v$, the graph $F^v$ computed in $G'$, contains $F^v$. The transitions at $v$ of the new system are
those of the old one except the transition of $R_p$, and their old linking paths are edge-disjoint paths
of $F^v\subseteq F'^v$.
\end{proof}

\begin{lemma}[base]\label{lem:base}
If $\Ev=\emptyset$, then $\emptyset$ is realisable, with $n/2$ paths.
\end{lemma}

\begin{proof}
All degrees are odd and positive, so Theorem~\ref{thm:lovasz} gives a partition with exactly one end at
each vertex, which is $b_\emptyset$.
\end{proof}

\begin{lemma}[extension]\label{lem:ext}
Let $M$ be a transversal of $P$, let $p=\{a,c\}$ with $a\in M$, and let $G'=G-E(R_p)$. In $G'$ both $a$ and
$c$ are odd, and $M-a$ is a transversal of $P-p$. The following are equivalent.
\begin{enumerate}
\item[(i)] $G'$ has a realisation of $M-a$ in which the path $Q_c$ ending at $c$ satisfies
$V(Q_c)\cap V(R_p)=\{c\}$.
\item[(ii)] $G$ has a realisation of $M$ in which $R_p$ is a terminal segment, ending at $a$, of some
path.
\end{enumerate}
\end{lemma}

\begin{proof}
(i)$\Rightarrow$(ii). $Q_c$ and $R_p$ share only $c$, which is an end of both, so $Q_c\cup R_p$ is a path.
Its ends are $a$ and the far end of $Q_c$. Replacing $Q_c$ by it removes the only end at $c$ and adds a
second end at $a$. On all other vertices the profiles of $M-a$ in $G'$ and of $M$ in $G$ agree.

(ii)$\Rightarrow$(i). Let $Q^\ast$ have $R_p$ as a terminal segment ending at $a$. Since $c\notin M$ has no
end, $Q^\ast$ continues beyond $c$. Deleting the edges of $R_p$ leaves a path $Q$ ending at $c$, with
$V(Q)\cap V(R_p)=\{c\}$. Now $c$ has one end and $a$ has one, which is the profile of $M-a$ in $G'$. Since
$b(c)=1$, $Q$ is the path $Q_c$.
\end{proof}

So route descent proves TRC for $(G,P)$ at $M$ if, for some $p$, the induction hypothesis on
$G-E(R_p)$ can be strengthened to deliver (i). The order of deletion, the choice of realisation, and
the route system producing a given pairing are all free.

\begin{proposition}[existential form and the Lean spine]\label{prop:spine}
Let $G$ be connected of even order and $y\in\Ev$. If some realisable $M\subseteq\Ev$ with $y\in M$ and
$|M|\le|\Ev|/2$ exists, then for every $k\le|\Ev|/2$ there is $R\subseteq\Ev$ with $|R|=k$, $y\notin R$,
and $\Ev\setminus R$ realisable. In particular, if every connected graph of even order has, for each of
its even vertices $y$, a realisable transversal of a Thm-21 pairing containing $y$, then Gallai's
conjecture holds.
\end{proposition}

\begin{proof}
$\Ev\setminus M$ has at least $|\Ev|/2\ge k$ elements and avoids $y$. Choose $R\subseteq\Ev\setminus M$
with $|R|=k$. Then $\Ev\setminus R\supseteq M$ is realisable by Lemma~\ref{lem:up}. This is the statement
$\mathrm{Desc}\,H\,k\,y\,\emptyset$ of the Lean library, so the hypothesis $\mathrm{BranchBase}_k$ holds
for every $k$, without its side conditions on $y$. The Lean theorem \texttt{gallai\_of\_branchBase}
gives Gallai's conjecture. The last claim applies this to the transversal containing $y$; Thm-21
systems exist by~\cite[Theorem~21]{EMM}.
\end{proof}

So for Gallai's conjecture we need only an \emph{existential} TRC. At each deletion we may choose
which end of $p$ to extend, except at the pair containing the root.

\section{Minimum $T$-joins}\label{sec:min}

A $T$-join of $G$ is \emph{minimum} if it has the fewest edges among all $T$-joins.

\begin{lemma}\label{lem:min}
Let $J$ be a minimum $T$-join of $G$ for $T=\Ev$, and let $\cR$ be any decomposition of $J$ into
edge-disjoint paths pairing $\Ev$.
\begin{enumerate}
\item $J$ is a forest. Hence Theorem~21 of~\cite{EMM} applies to $J$, and $J$ carries a Thm-21 system.
\item For every cycle $C$ of $G$, $|E(C)\cap J|\le|E(C)\setminus J|$.
\item Every segment $R[x,y]$ of every route $R\in\cR$ is a shortest $x$--$y$ path of $G$. In particular
every route is an induced path of $G$.
\item If $p=\{a,c\}$ and $G'=G-E(R_p)$, then $J-E(R_p)$ is a minimum $T'$-join of $G'$, where
$T'=\Ev(G')$.
\item A $T$-join that is a matching is minimum.
\end{enumerate}
\end{lemma}

\begin{proof}
For any even subgraph (edge set with all degrees even) $E$, the symmetric difference $J\triangle E$ is again
a $T$-join, of size $|J|-|J\cap E|+|E\setminus J|$. Minimality gives $|J\cap E|\le|E\setminus J|$.
Taking $E=E(C)$ gives (2), and if $C\subseteq J$ it gives $|E(C)|\le0$, so (1). For (3), let $Q$ be an
$x$--$y$ path and $E=E(R[x,y])\triangle E(Q)$, an even subgraph. Then
$|J\cap E|\ge|E(R[x,y])\setminus E(Q)|$ and $|E\setminus J|\le|E(Q)\setminus E(R[x,y])|$. So
$|E(R[x,y])\setminus E(Q)|\le|E(Q)\setminus E(R[x,y])|$, that is, $|R[x,y]|\le|Q|$. A chord of a route
would give a shorter path between two of its vertices. For (4), if $K'$ were a $T'$-join of $G'$ with
$|K'|<|J|-|R_p|$, then $K'\cup E(R_p)$, a disjoint union, would be a $T$-join of $G$: the parities of
$K'$ and $R_p$ add up to $T'\triangle\{a,c\}=T$. It would have fewer edges than $J$. For (5), every
$T$-join has at least $|T|/2$ edges, since every vertex of $T$ needs an edge and an edge covers at most
two vertices.
\end{proof}

By (1) and (4) and Lemma~\ref{lem:hered}, the class of \emph{minimum Thm-21 systems}, those built on a
minimum $T$-join, is closed under route deletion, and route descent can stay inside it.

\begin{corollary}[no trivial obstruction]\label{cor:triv}
In a route system on a minimum $T$-join, let $p=\{a,c\}$ and $G'=G-E(R_p)$. Then $c$ has a neighbour in
$G'$ outside $V(R_p)$, and so does $a$.
\end{corollary}

\begin{proof}
$c$ has degree at least two and exactly one edge on $R_p$. Any other neighbour of $c$ on $R_p$ would give a
chord of $R_p$, against Lemma~\ref{lem:min}(3).
\end{proof}

\begin{remark}\label{rem:triv}
With spanning-tree $T$-joins the extension (i) of Lemma~\ref{lem:ext} fails in a few percent of cases.
Most failures are of the trivial kind excluded by Corollary~\ref{cor:triv}: $c$ has degree two and the
route runs through its other neighbour, so $Q_c$ cannot leave the route. The only instances in which no
route at all admits the extension were all of this kind. With minimum $T$-joins no failure of (i) was
found, for any route or any transversal (Section~\ref{sec:conj}).
\end{remark}

\section{Route-respecting realisations}\label{sec:rr}

\begin{definition}
Let $M$ be a transversal of $P$. For $p\in P$ write $a_p\in M$ and $c_p\notin M$ for its two vertices
(head and tail). A realisation of $M$ is \emph{route-respecting} if every route $R_p$ is a terminal
segment, ending at $a_p$, of some path. A Lov\'asz decomposition $\cD$ of the free graph $F$ (which has all
degrees odd, by Lemma~\ref{lem:basic}) is \emph{$M$-compatible} if, for every path $Q\in\cD$:
\begin{enumerate}
\item[(RR1)] if an end $s$ of $Q$ is a tail $c_p$, then $V(Q)\cap V(R_p)=\{s\}$;
\item[(RR2)] if both ends of $Q$ are tails $c_p,c_q$, then $V(R_p)\cap V(R_q)=\emptyset$.
\end{enumerate}
\end{definition}

\begin{proposition}\label{prop:rr}
Route-respecting realisations of $M$ correspond bijectively to $M$-compatible Lov\'asz decompositions of
$F$: append every route at its tail. In particular, if $F$ has an $M$-compatible Lov\'asz decomposition,
then $M$ is realisable.
\end{proposition}

\begin{proof}
Let $\cD$ be $M$-compatible. For $Q\in\cD$ attach $R_p$ at every end of $Q$ that is a tail $c_p$. By
(RR1) and (RR2) the result is a path: the attached routes meet $Q$ only at its ends, and two attached
routes are disjoint. Every edge of $J$ is used exactly once, since every route has exactly one tail and
that tail is the end of exactly one path of $\cD$. As for ends: a tail loses its unique end, so it has
none; a head $a_p$ keeps its end in $\cD$ and gains the end of $R_p$, so it has two; every odd vertex of
$G$ keeps one. This is $b_M$. Conversely, in a route-respecting realisation each path contains at most two
routes as terminal segments, one at each end. Every edge of $J$ lies in such a segment, so deleting all
routes leaves only edges of $F$. A path does not end at a tail, so each path properly extends its
terminal routes, and deleting them leaves a nonempty path. The pieces have one end at every vertex and
satisfy (RR1) and (RR2), since they came from paths.
\end{proof}

By Proposition~\ref{prop:rr}, TRC for minimum Thm-21 systems would follow from the following
conjecture, which concerns only graphs with all degrees odd.

\begin{conjecture}[RR]\label{conj:rr}
For every minimum Thm-21 system and every transversal $M$, the free graph has an $M$-compatible
Lov\'asz decomposition.
\end{conjecture}

The decomposition must depend on $M$. There are minimum Thm-21 systems, already among the order-$12$
counterexamples to EMM, for which no single Lov\'asz decomposition of $F$ is compatible with all
transversals at once (Section~\ref{sec:conj}). The linked-transition condition cannot be dropped either:
for minimum $T$-joins decomposed into routes by \emph{arbitrary} transitions, some transversals are not
even realisable (Section~\ref{sec:conj}).

When every route is a single edge, $J=N$ is a perfect matching of $G[\Ev]$. It is minimum by
Lemma~\ref{lem:min}(5), and every vertex of $\Ev$ has $J$-degree one, so the transitions are linked
trivially: the end edge $vz$ is linked in $F^v$, whose only odd vertices are $v$ and $z$. Condition (RR2)
is automatic, because distinct edges of a matching are disjoint. So Conjecture~\ref{conj:rr} becomes:

\begin{conjecture}[$L_N$]\label{conj:LN}
Let $F$ have all degrees odd, and let $c_1a_1,\dots,c_ma_m$ be pairwise disjoint pairs of non-adjacent
vertices of $F$. Then $F$ has a Lov\'asz decomposition in which, for every $i$, the path ending at $c_i$
does not contain $a_i$.
\end{conjecture}

\begin{proposition}\label{prop:LN}
Conjecture~\ref{conj:LN} implies Gallai's conjecture for every connected graph $G$ of even order whose
even vertices induce a subgraph with a perfect matching. Moreover, it implies that every orientation of
such a matching is realisable.
\end{proposition}

\begin{proof}
Let $N$ be a perfect matching of $G[\Ev]$, oriented $c_i\to a_i$, and let $F=G-N$. Every vertex has odd
degree in $F$, and the pairs $c_ia_i$ are disjoint and non-adjacent in $F$. Apply Conjecture~\ref{conj:LN}
and Proposition~\ref{prop:rr} with the single-edge routes $c_ia_i$. The realisation of $\{a_i\}$ has
$|\Od|/2+|\Ev|/2=n/2$ paths.
\end{proof}

This class contains every Eulerian graph with a perfect matching. For such graphs Gallai's conjecture is
open. The case $m=1$ of Conjecture~\ref{conj:LN} is:

\begin{conjecture}[$L_1$]\label{conj:L1}
If $F$ has all degrees odd and $c,a$ are non-adjacent, then some Lov\'asz decomposition of $F$ has
$a\notin V(Q_c)$.
\end{conjecture}

The next section proves $L_1$ when $c$ and $a$ have no common neighbour.

\section{Lov\'asz decompositions that separate vertices}\label{sec:sep}

Say that a path partition \emph{separates} $c$ and $a$ if no path contains both. A set $X$ of
vertices is \emph{packed} if its vertices are pairwise at distance at least three, that is, no two are
adjacent and no two have a common neighbour.

\begin{theorem}[separation]\label{thm:sep}
Let $K$ have at most one vertex of positive even degree, and let $X$ be a packed set of vertices. Then
$K$ has a Lov\'asz decomposition in which every path contains at most one vertex of $X$. In particular,
if $\mathrm{dist}_K(c,a)\ge3$, then some Lov\'asz decomposition of $K$ separates $c$ and $a$.
\end{theorem}

\begin{proof}
Induction on $\nu(K)=2|E(K)|+[\text{$K$ has no vertex of positive even degree}]$. Call a partition
\emph{good} if each of its paths contains at most one vertex of $X$. Deleting edges and adding pendant
edges at new vertices keeps $X$ packed, since a new pendant vertex lies on no path between two old
vertices. The key observation is this. If $\cD$ is good and $\cD^+$ is obtained by Lemma~\ref{lem:fan}
with target $u\notin X$, then $\cD^+$ is good. Indeed, all vertex sets are preserved except one, which
gains $u\notin X$.

\emph{Case I: $K$ has a vertex $z$ of positive even degree.} Then $z$ has a neighbour $u\notin X$. If
$z\in X$, no neighbour of $z$ lies in $X$. If $z\notin X$, at most one of its (at least two)
neighbours lies in $X$, since two would have the common neighbour $z$. Let $K'=K-uz$. In $K'$ the
vertex $z$ is odd. The vertex $u$ was odd, since $z$ is the only even vertex, so it is now even or
isolated. All other degrees are unchanged. So $K'$ has at most one vertex of positive even degree, and
$\nu(K')<\nu(K)$. By induction $K'$ has a good Lov\'asz decomposition $\cD'$. Then
$b_{\cD'}(z)=1$ and $b_{\cD'}(u)=0$. Every neighbour of $u$ in $K'$ has odd degree, hence one end, so
$Z_{\cD'}(u)=\emptyset$. Lemma~\ref{lem:fan} with $t=z$ and target $u$ gives a decomposition of $K$
with ends $b_{\cD'}-\mathbf 1_z+\mathbf 1_u=d_K\bmod2$. It is good, as $u\notin X$.

\emph{Case II: no vertex has positive even degree, and some $x\notin X$ has degree at least $3$.}
Choose any $u\in N(x)$, and let $H$ be $K-ux$ together with a new vertex $\ell$ and the edge $u\ell$. In $H$
the vertex $x$ has positive even degree $d(x)-1$, and every other vertex has odd degree or degree zero.
$H$ has as many edges as $K$ but has an even vertex, so $\nu(H)<\nu(K)$. By induction $H$ has a good
Lov\'asz decomposition $\cD_H$. The leaf $\ell$ is an end of its path $Q_\ell$. Delete $\ell$. If
$Q_\ell=u\ell$, delete this path; then $u$ loses its only end. Otherwise $u$ was interior to $Q_\ell$ and
becomes an end of $Q_\ell-\ell$, while keeping its own end; so it has two. The result $\cD'$ is a good
partition of $K'=K-ux$. It has $b_{\cD'}(x)=0$, $b_{\cD'}(u)\in\{0,2\}$, and
$b_{\cD'}(w)=d_K(w)\bmod2$ for $w\ne u,x$.
\begin{itemize}
\item If $b_{\cD'}(u)=0$, add the one-edge path $ux$. It contains at most one vertex of $X$, as
$x\notin X$.
\item If $b_{\cD'}(u)=2$, use Lemma~\ref{lem:fan} with $t=u$ and target $x$. Every neighbour of $x$ in
$K'$ has odd degree in $K'$ and one end; $u$ itself is not a neighbour of $x$ in $K'$. So
$Z_{\cD'}(x)=\emptyset<2$. The result has one end at $u$ and at $x$, and is good as $x\notin X$.
\end{itemize}

\emph{Case III: no vertex has positive even degree, every vertex outside $X$ has degree at most two,
and some $v\notin X$ has positive degree.} Then $d(v)=1$. Let $w$ be its neighbour and $K'=K-vw$. In
$K'$ the vertex $v$ is isolated. The vertex $w$ was odd, so it is now even or isolated. All other
degrees are unchanged, and $\nu(K')<\nu(K)$. A good Lov\'asz decomposition of $K'$ has no end at $w$.
Adding the one-edge path $vw$ gives a Lov\'asz decomposition of $K$. It is good, since $v\notin X$.

\emph{Case IV: every vertex outside $X$ has degree zero.} An edge would have an end outside $X$, since
no two vertices of $X$ are adjacent; so $K$ has no edges, and the empty partition is good.

These cases are exhaustive. If there is no even vertex and no vertex outside $X$ of degree at least
three, then the vertices outside $X$ have degree at most one, being odd or isolated; so Case~III or
Case~IV applies. This completes the induction. For the last statement, take $X=\{c,a\}$.
\end{proof}

Theorem~\ref{thm:sep} is formalised in Lean as \texttt{separation\_set}, with the case
$X=\{c,a\}$ as \texttt{separation} (Section~\ref{sec:lean}).

\begin{corollary}[$L_1$ at distance three]\label{cor:L1}
If all degrees of $F$ are odd and $c,a$ are non-adjacent with no common neighbour, then $F$ has a
Lov\'asz decomposition in which $Q_c$ avoids $a$ and $Q_a$ avoids $c$.
\end{corollary}

\begin{corollary}\label{cor:onepair}
Let $G$ have exactly two even vertices $a$ and $c$, and suppose $ac\in E(G)$ lies in no triangle. Then
$\{a\}$ has a realisation in which $ca$ is a terminal edge at $a$, and $\{c\}$ has one in which $ac$ is a
terminal edge at $c$.
\end{corollary}

\begin{proof}
In $F=G-ac$ all degrees are odd, and $c,a$ are non-adjacent with no common neighbour. Apply
Corollary~\ref{cor:L1} and Proposition~\ref{prop:rr} with the single route $ca$.
\end{proof}

\begin{corollary}\label{cor:packedroutes}
Let $(P,\cR)$ be a route system for $G$ with union $J$, and suppose $V(J)$ is packed in the free graph
$F=G-J$. Then one Lov\'asz decomposition of $F$ is $M$-compatible for every transversal $M$; so all
transversals are realised by appending the routes to the paths of a single partition of $F$.
\end{corollary}

\begin{proof}
By Theorem~\ref{thm:sep}, some Lov\'asz decomposition of $F$ has every path meeting $V(J)$ at most once.
Then (RR1) holds for every choice of tails, and (RR2) holds vacuously, since no path has two ends in
$V(J)$.
\end{proof}

The hypothesis of Corollary~\ref{cor:packedroutes} forces $G[V(J)]=J$, a forest, so $G[\Ev]$ is a forest.
These graphs lie in Pyber's class, and the corollary adds only the uniform realisation of all
transversals.

\begin{example}[distance two]\label{ex:dist2}
Separation can fail at distance two in the presence of an even vertex. Let $K$ have vertices
$z,u_1,u_2,c,a,\ell_c,\ell_a$ and edges $zu_1,zu_2$, $u_ic$, $u_ia$ ($i=1,2$), $c\ell_c$, $a\ell_a$.
Then $z$ is the only even vertex, and $\mathrm{dist}(c,a)=2$. Suppose a Lov\'asz decomposition
separates $c$ and $a$. The vertex $z$ has no end, so some path $Q$ contains $u_1zu_2$. Each $u_i$ has
degree three and exactly one end, hence exactly one passage. If $u_i$ were an end of $Q$, its end edge would
be $u_iz$ and its passage would be $\{c,a\}$, and the path through that passage would contain both $c$ and
$a$. So $Q$ continues beyond $u_1$ to some $w_1\in\{c,a\}$ and beyond $u_2$ to some $w_2\in\{c,a\}$.
If $w_1=w_2$, then $Q$ repeats a vertex; otherwise $Q$ contains both $c$ and $a$. Either way we have a
contradiction, so no Lov\'asz decomposition of $K$ separates $c$ and $a$. Yet $L_1$ holds here. The
paths $c\,u_1$, $\ \ell_c\,c\,u_2\,z\,u_1\,a$ and $u_2\,a\,\ell_a$ partition the eight edges, have one end
at each odd vertex and none at $z$, and the path ending at $c$ is $c\,u_1$, which avoids $a$. Adding a
leaf at $z$ makes all degrees odd, and then a separating decomposition exists again (found by computation;
not used).
\end{example}

\begin{remark}[how far Lov\'asz's induction proves $L_1$]\label{rem:L1res}
For $L_1$ itself the invariant is weaker: only the path \emph{ending} at $c$ must avoid $a$. In
Lemma~\ref{lem:fan} with target $g\neq a$ the chain vertices other than $t$ are neighbours of $g$. If
$c\notin N(g)\cup\{t\}$, the path ending at $c$ keeps its vertex set, or gains $g$. If $c$ is a chain
vertex, the end at $c$ moves onto the path containing the edge $cg$. So the invariant survives whenever
$g\ne a$ and either $c\notin N(g)\cup\{t\}$ or the path through $cg$ avoids $a$. Running the induction of
Theorem~\ref{thm:sep} with this rule leaves two residual configurations:
\begin{itemize}
\item the even vertex $z$ has $N(z)\subseteq N[c]\cup\{a\}$;
\item all degrees are odd, the non-isolated vertices are $c$, $N(c)$ and $a$, $N(a)\subseteq N(c)$, and
every vertex of $N(c)$ has degree at least three.
\end{itemize}
Leaves adjacent to $c$ finish at once, with the one-edge path $Q_c$. Leaves outside $N[c]$ are deleted and
restored. Example~\ref{ex:dist2} shows that the residual cases need an invariant strictly between $L_1$ and
separation.
\end{remark}

\section{Classes reachable by Fan's lemma}\label{sec:classes}

\begin{theorem}[Fan-isolated routes]\label{thm:isolated}
Let $(P,\cR)$ be a route system for $G$ such that every vertex $x$ of every route $R_p$ satisfies
\[
N_G(x)\cap\Ev\subseteq N_{R_p}(x).
\]
Then every transversal of $P$ is realisable.
\end{theorem}

\begin{proof}
Induction on $|P|$, with Lemma~\ref{lem:base} as the base. The hypothesis passes to
$G'=G-E(R_p)$ with the remaining routes, because $N_{G'}(x)\cap\Ev(G')\subseteq N_G(x)\cap\Ev$
(Lemma~\ref{lem:hered}). Let $M$ be a transversal, $p=\{a,c\}$ with $a\in M$, and write
$R_p=r_0r_1\cdots r_t$ with $r_0=c$ and $r_t=a$. By induction there is a realisation $\cD_0$ of $M-a$ in
$G'$. Put the edges $r_{i-1}r_i$ back one at a time, for $i=1,\dots,t$. At step $i$, use
Lemma~\ref{lem:fan} in $G_{i-1}=G'+r_0r_1+\dots+r_{i-2}r_{i-1}$ with $t=r_{i-1}$ and target $r_i$. This
moves one end from $r_{i-1}$ to $r_i$. So after step $i-1$ the ends are $b_{i-1}=b_0-\mathbf
1_c+\mathbf 1_{r_{i-1}}$ (with $b_0$ the profile of $M-a$ in $G'$).
\begin{itemize}
\item $b_{i-1}(r_{i-1})\ge1$: for $i=1$ because $c$ is odd in $G'$, and for $i\ge2$ because $r_{i-1}$
received an end.
\item A vertex with no end in $\cD_{i-1}$ is an unmarked vertex of $\Ev$ (here $c$ counts, for $i\ge2$).
So $Z(r_i)\subseteq N_{G_{i-1}}(r_i)\cap\Ev\subseteq N_{R_p}(r_i)$. But neither $R_p$-neighbour of
$r_i$ is adjacent to $r_i$ in $G_{i-1}$: the edges $r_{i-1}r_i$ and $r_ir_{i+1}$ have not yet been added.
So $Z(r_i)=\emptyset$.
\end{itemize}
Thus each transfer is legal, and after step $t$ the ends are $b_0-\mathbf 1_c+\mathbf 1_a$, which is the
profile of $M$ in $G$.
\end{proof}

The hypothesis forces $G[\Ev]$ to be a matching. Let $e\in\Ev$ be the end of $R_p$. If $e$ is also
interior to a route $R_q$, then its even neighbours lie in $N_{R_p}(e)\cap N_{R_q}(e)$, which is empty
because the two routes are edge-disjoint. Otherwise its only possible even neighbour is its neighbour
$y$ on $R_p$. If $y$ is not the other end of $R_p$, then $y$ is interior to $R_p$ and the end of another
route $R_q$, and the hypothesis at $y$ forces $e$ to be its neighbour on $R_q$; then the edge $ey$ lies on
two routes, which is impossible. So every even vertex has at most one even neighbour, namely its partner
along a one-edge route. Hence these graphs lie in Pyber's class, and Theorem~\ref{thm:isolated} is new
only as a statement about \emph{all} transversals.

The same argument works along any order of the routes, provided each route is clean at the moment it
is put back. This is the mechanism behind the existential form of route descent.

\begin{definition}[route schedule]
Let $(P,\cR)$ be a route system with free graph $F$. A \emph{route schedule} is an ordering
$p_1,\dots,p_m$ of $P$ together with a choice of a tail $c_j$ and a head $a_j$ in each $p_j$, with the
following property. Put $G_0=F$ and $G_j=G_{j-1}\cup R_{p_j}$, and let $U_j=\{c_1,\dots,c_j\}$. Then for
every $j$, no vertex $x\neq c_j$ of $R_{p_j}$ has a neighbour in $G_{j-1}$ that lies in $U_{j-1}\cup\{c_j\}$.
\end{definition}

\begin{theorem}[route schedules]\label{thm:schedule}
If $(P,\cR)$ has a route schedule, then $\{a_1,\dots,a_m\}$ is realisable.
\end{theorem}

\begin{proof}
We show by induction on $j$ that $G_j$ has a path partition with two ends at $a_1,\dots,a_j$, none at
$c_1,\dots,c_j$, and $d_{G_j}(v)\bmod2$ ends at every other vertex. For $j=0$ this is
Theorem~\ref{thm:lovasz}, since $F$ has all degrees odd (Lemma~\ref{lem:basic}). Every vertex of a later
pair is odd in $G_j$, because each route through it contributes two edges. For the step, write
$R_{p_j}=r_0r_1\cdots r_t$ with $r_0=c_j$ and $r_t=a_j$, and add its edges one at a time by
Lemma~\ref{lem:fan}, moving an end from $r_{i-1}$ to $r_i$, exactly as in the proof of
Theorem~\ref{thm:isolated}. At step $i$ the vertex $r_{i-1}$ has at least one end. The vertices without
ends are those of $U_{j-1}$, together with $c_j$ when $i\ge2$. The neighbours of $r_i$ in the current
graph are its neighbours in $G_{j-1}$, since its two edges on $R_{p_j}$ are not yet present. By the
schedule property none of them has zero ends, so the transfer is legal.
\end{proof}

Theorem~\ref{thm:isolated} is the case in which every order and every orientation is a schedule. For
one-edge routes the condition says that each head $a_j$ is not adjacent to the earlier tails
$c_1,\dots,c_{j-1}$. These are exactly the Fan schedules (semi-induced matchings) of the
endpoint-structure manuscript. For minimum $T$-joins the condition at $c_j$ itself is automatic, since
routes have no chords (Lemma~\ref{lem:min}(3)), so it reduces to the earlier tails. A schedule whose
heads include a prescribed vertex $y$ gives the existential form at $y$
(Proposition~\ref{prop:spine}).

\begin{corollary}\label{cor:upm}
If $G[\Ev]$ has a unique perfect matching $N$, then every orientation of $N$ (choice of one vertex of
each edge) is realisable.
\end{corollary}

\begin{proof}
Fix an orientation $c_i\to a_i$. Call $i\to j$ if $a_i$ is adjacent to $c_j$ with $j\neq i$. A directed
cycle $i_1\to\cdots\to i_r\to i_1$ would give the cycle $c_{i_1}a_{i_1}c_{i_2}a_{i_2}\cdots$ in $G[\Ev]$,
alternating between edges of $N$ and edges outside $N$. Switching $N$ along it gives a second perfect
matching. So the relation is acyclic, and some $i$ has no successor: $a_i$ has no even neighbour other
than $c_i$ among the unmarked vertices $c_j$. Delete the edge $c_ia_i$. In $G'=G-c_ia_i$ the induced
subgraph on $\Ev(G')=\Ev\setminus\{a_i,c_i\}$ is $G[\Ev]-\{a_i,c_i\}$, whose perfect matching
$N-c_ia_i$ is again unique. By induction $\{a_j:j\ne i\}$ is realisable in $G'$. In that realisation the
zero-end vertices are the $c_j$ with $j\neq i$, none of which is adjacent to $a_i$, while $c_i$ has one
end. Lemma~\ref{lem:fan} with $t=c_i$ and target $a_i$ restores $c_ia_i$ and moves the end of $c_i$ to $a_i$.
The base case $N=\emptyset$ is Lemma~\ref{lem:base}.
\end{proof}

Corollary~\ref{cor:upm} is the acyclic-ordering form of the Fan schedules of the endpoint-structure
manuscript (\texttt{realisable\_of\_fanSchedule}). Uniqueness of the matching is exactly the condition
that every orientation admits such an order. Gallai's conjecture for these graphs already follows from
those schedules.

\section{Conjectures and their status}\label{sec:conj}

The following computations guided the statements above; no proof uses them. An exact oracle extends the
depth-first search of the project's realisability oracle to constrained queries. It was validated
against an independent enumeration of all path partitions of small graphs. All realisations it reports
are re-checked by an independent partition checker.

\begin{itemize}
\item \emph{TRC.} No counterexample. All transversals of $3{,}600$ Thm-21 pairings of the twelve order-$12$
counterexamples to EMM ($86{,}400$ transversals) are realisable. So are the transversals of
random graphs of orders $10$ and $12$ (about $60{,}000$ transversals), all $1{,}280$ transversals of $20$
minimum Thm-21 pairings of the order-$18$ graph of minimum degree three~\cite[Thm~41]{EMM}, and $128$
sampled transversals of the order-$28$ graph built from the Petersen graph in~\cite{EMM}. All
$1{,}280$ transversals of $20$ Thm-21 pairings of the order-$16$ even-degree hardness graph built from the
prism~\cite[Thm~45]{EMM} are realisable; on its Petersen analogue (order $28$), $115$ of $128$ sampled
transversals were found realisable and $13$ were undecided within the search budget. In these graphs the
marks automatically fall inside the forcing gadgets, as they must for transversals.
\item \emph{The extension step.} With spanning-tree $T$-joins, (i) of Lemma~\ref{lem:ext} fails for
$4$--$5\%$ of the pairs $(M,p)$ in sparse random graphs; $95\%$ of the failures are trivial. With
\emph{minimum} $T$-joins it never failed, in about $220{,}000$ pairs $(M,p)$, including $8{,}896$ on
$80$ random graphs of order $14$.
\item \emph{RR.} An $M$-compatible decomposition was found for every transversal tested with minimum
$T$-joins ($1{,}276$ of order $10$, $1{,}280$ on the order-$18$ graph, $128$ on the order-$28$ graph). A
single decomposition compatible with every transversal at once fails to exist in two of the twelve
order-$12$ graphs, even with minimum $T$-joins.
\item \emph{Linking is needed.} For minimum $T$-joins decomposed into routes by random
transitions, without the linking condition, $4$ of $1{,}834$ transversals in random graphs of order $10$
were not realisable.
\item \emph{$L_1$, $L_N$.} No failure of $L_1$ in $70{,}978$ non-adjacent pairs (orders $\le10$), nor of
the stronger statement that $Q_c\subseteq N[c]$ can be achieved ($19{,}040$ vertices). No failure of
$L_N$ in $3{,}732$ instances of orders up to $14$ with up to seven pairs. When the pairs form a perfect matching $N$ of $V(F)$ (the Eulerian case: $F+N$ has all
degrees even), no failure in $2{,}332$ decided
instances of orders $8$--$20$ ($157$ more were undecided within the search budget). In $2{,}421$ of the
$2{,}489$ instances some pair is at distance two, where the symmetric version below fails.
\item \emph{Separation.} The symmetric version of $L_N$ (no path contains both vertices of any pair) is
false when some pair is at distance two, failing in about a fifth of random instances. It held in all
$2{,}892$ tested instances of orders $8$--$16$ in which all pairs are at distance at least three,
including $925$ perfect matchings of the vertex set ($102$ further instances were undecided within the
search budget).
\end{itemize}

\begin{conjecture}[multi-pair separation]\label{conj:msep}
If $K$ has all degrees odd and $\{s_1,t_1\},\dots,\{s_m,t_m\}$ are pairwise disjoint pairs at distance at
least three, then some Lov\'asz decomposition of $K$ separates every pair.
\end{conjecture}

Theorem~\ref{thm:sep} is the case of Conjecture~\ref{conj:msep} in which the union of the pairs is
packed. Conjecture~\ref{conj:msep} implies Conjecture~\ref{conj:LN} for pairs at distance at least three. So it
would give Gallai's conjecture, with every orientation realisable, for every graph whose even vertices
induce a subgraph having a perfect matching none of whose edges lies in a triangle of $G$. This includes
every triangle-free Eulerian graph with a perfect matching. The proof of Theorem~\ref{thm:sep} does not
extend directly: when every vertex lies in a pair, as for an Eulerian graph, every Fan target is
paired, and the path that gains the target may contain its partner.

\section{Comparison with known results}\label{sec:compare}

\begin{itemize}
\item Lov\'asz's theorem~\cite{L68} is the base of route descent. Theorem~\ref{thm:sep} strengthens its
strong form: the decomposition can be chosen so that every path meets a prescribed packed set at most
once. We have not found this in the literature.
\item Pyber~\cite{Pyb96} (even subgraph a forest), Fan~\cite{Fan05} (blocks triangle-free of maximum
degree at most three) and Botler and Sambinelli~\cite{BS21} prove Gallai's conjecture under structural
conditions on $G[\Ev]$. Theorem~\ref{thm:isolated} lies inside Pyber's class. Corollary~\ref{cor:upm}
follows from the Fan schedules of the endpoint-structure manuscript. Both are new only as statements that
\emph{every} transversal is realisable.
\item Conjectures~\ref{conj:LN} and~\ref{conj:msep} would give new classes: graphs whose even subgraph has
a perfect matching (respectively one avoiding triangles). These are not covered by the results above,
because the even subgraph may have large degrees and many cycles.
\item The heredity of minimum $T$-joins (Lemma~\ref{lem:min}(4)) is elementary. Its use here is to make
route descent avoid the only obstruction ever observed.
\end{itemize}

\section{Formalisation}\label{sec:lean}

The following are proved in Lean~4 with Mathlib, in the project library \texttt{lean/GallaiLean}. The
files are sorry-free, and each theorem listed depends only on the standard axioms (\texttt{propext},
\texttt{Classical.choice}, \texttt{Quot.sound}).
\begin{itemize}
\item \texttt{RouteDescent.lean}.
  \begin{itemize}
  \item \texttt{gallai\_of\_pairingTRC}: if every connected graph of even order has a pairing of its
  even vertices all of whose transversals are realisable, Gallai's conjecture holds. The variant
  \texttt{gallai\_of\_pairingTRCExists} needs only one realisable transversal through each even vertex.
  Both go through \texttt{branchBase\_of\_rootedHalfSets}, which is Proposition~\ref{prop:spine}.
  \item \texttt{IsEvenPairing.of\_isPairing}: the pairings of the formalised cut-selection theorem
  (\texttt{cs\_all}) are pairings in this sense, with the same transversals.
  \item \texttt{isMinTJoin\_sdiff} is Lemma~\ref{lem:min}(4), and
  \texttt{isMinTJoin\_of\_two\_mul\_card\_le} is Lemma~\ref{lem:min}(5).
  \item \texttt{realisable\_of\_even\_matching} is Theorem~\ref{thm:isolated} when every route is a
  single edge, that is, when $G[\Ev]$ is a perfect matching.
  \item \texttt{PathDecomp.route\_append} is Lemma~\ref{lem:ext}, (i)$\Rightarrow$(ii), the inductive
  step of route descent: for a path $R$ of $G$ from $c$ to $a$ and a decomposition of $G-E(R)$ whose path
  at $c$ meets $R$ only in $c$, it gives a decomposition of $G$ of the same size whose end counts are
  those of the old one with one end moved from $c$ to $a$.
  \end{itemize}
\item \texttt{FanSupport.lean}: Lemma~\ref{lem:fan} (\texttt{fan\_single\_supp},
\texttt{fan\_grouped\_supp}, \texttt{fan\_exec\_supp}).
\item \texttt{Separation.lean}: Theorem~\ref{thm:sep} (\texttt{separation\_set}, and
\texttt{separation} for $X=\{c,a\}$) and Corollary~\ref{cor:L1} (\texttt{lovasz\_avoid}). The formal
proof follows Cases I--IV above. The leaf of Case~II lives on the type \texttt{Option W} and is removed
by a projection lemma; a strengthened copy, \texttt{Retract.delete\_supp}, records the vertex sets.
\end{itemize}
Not formalised: Lemma~\ref{lem:hered}, Lemma~\ref{lem:ext} (ii)$\Rightarrow$(i), Lemma~\ref{lem:min}(1)--(3),
Proposition~\ref{prop:rr}, Theorem~\ref{thm:isolated} for routes of length at least two, and
Corollary~\ref{cor:upm}, which is close to the existing \texttt{realisable\_of\_fanSchedule}.

\section{Open problems}\label{sec:open}
\begin{enumerate}
\item Prove $L_1$ (Conjecture~\ref{conj:L1}) at distance two. The separation invariant of
Theorem~\ref{thm:sep} is too strong there (Example~\ref{ex:dist2}), so a weaker invariant is needed. It
must still survive Fan transfers whose chains pass through $c$.
\item Prove Conjecture~\ref{conj:msep}, or Conjecture~\ref{conj:LN}, for two pairs.
\item Is Conjecture~\ref{conj:rr} true for routes of length two, that is, for star-forest $T$-joins?
\item Does every \emph{minimum} Thm-21 system admit route descent in every order of deletion, as the
computations suggest?
\end{enumerate}

\begin{thebibliography}{99}
\bibitem{BP19} M.~Bonamy and T.~J.~Perrett, Gallai's path decomposition conjecture for graphs of small
maximum degree, Discrete Math.\ 342 (2019) 1293--1299.
\bibitem{BS21} F.~Botler and M.~Sambinelli, Towards Gallai's path decomposition conjecture, J.~Graph
Theory 97 (2021) 161--184.
\bibitem{EMM} J.~Williams, Prescribed ends in path decompositions: a cut inequality, cut selection, and
hardness, manuscript (\texttt{manuscripts/endpoint-min-max}), 2026; together with the project's
\texttt{research/synthesis.md}.
\bibitem{Fan05} G.~Fan, Path decompositions and Gallai's conjecture, J.~Combin.\ Theory Ser.~B 93 (2005)
117--125.
\bibitem{L68} L.~Lov\'asz, On covering of graphs, in: Theory of Graphs (Proc.\ Colloq.\ Tihany, 1966),
Academic Press, New York, 1968, 231--236.
\bibitem{Pyb96} L.~Pyber, Covering the edges of a connected graph by paths, J.~Combin.\ Theory Ser.~B 66
(1996) 152--159.
\end{thebibliography}
\end{document}
